\chapter{Teaching with the PRG32 Construction Kit}
\label{ch:kitteach}

This chapter is for the adult in the room. It explains what the PRG32
Construction Kit is, how to install and run it for a class, how the three
pupil chapters before it are meant to be taught, how to assess what children
learn, and how the Blocks track connects to the C and assembly tracks that
older students follow on the same platform. No programming experience is
assumed, although a teacher who has read Chapter~\ref{ch:platform} will
recognise every idea.

\section{What the Kit is, and what it is not}

The PRG32 Construction Kit is a small web application: a server written in
Python that a teacher runs on one computer, and a page that pupils open in a
browser \citep{prg32kit}. Pupils build programs from blocks, in the manner made
familiar by Scratch \citep{resnick2009}, using an editor built on the Blockly
library \citep{blockly}. Five components matter in the classroom.

\begin{center}
\begin{tabularx}{\linewidth}{@{}lX@{}}
\toprule
Component & What it gives the class \\
\midrule
Blocks editor & a 20-block vocabulary in five drawers: Game, State, Input \&
  Logic, Drawing, Audio \\
Simulator & a $320\times200$ game screen in the browser with an on-screen game
  pad, Play/Pause, single-frame \emph{Step}, Reset, and a live read-out of the
  variables \\
Sprite Lab & a pixel editor that saves sprites and converts them to C arrays \\
C generator & deterministic translation of a Blocks project into real \prg C
  source, shown beside the blocks \\
Packager and publisher & builds \fname{.prg32} cartridges when the \prg
  toolchain is present and submits bundles to a Cartridge Store \\
\bottomrule
\end{tabularx}
\captionof{table}{The parts of the Construction Kit \citep{prg32kit}.}
\end{center}

The design decision that distinguishes the Kit from a general-purpose blocks
environment is that it is \emph{not a toy dialect}. The three top-level blocks
are the \code{init}, \code{update}, and \code{draw} functions of
Chapter~\ref{ch:platform}; the drawing blocks are the calls of
Table~\ref{tab:gfxcore}; and the generated C compiles, unmodified, into a
cartridge that runs on the ESP32\nobreakdash-C6. A child of eight and a
first-year undergraduate are therefore programming the same machine through
the same interface, at different altitudes.

\begin{prgconcept}
The Kit's pipeline is: blocks $\rightarrow$ an intermediate description of the
game $\rightarrow$ the browser simulator \emph{and} generated C $\rightarrow$ a
\fname{.prg32} cartridge $\rightarrow$ the Cartridge Store. Every stage is
visible to the pupil, and none is magic.
\end{prgconcept}

It is equally important to know the limits. The simulator approximates the
runtime: sound is a browser tone, and random sequences differ from the
board's. The Blocks vocabulary covers one controller, integer variables,
rectangle collision, rectangles, pixels, text, an indexed palette, random
numbers, and simple notes. Tile maps, scrolling playfields, sprites, tracker
music, scores, and multiplayer exist in the cartridge ABI but have no blocks
yet; they are reached through the Kit's \emph{Advanced~C} editor, which builds
real cartridges but is not simulated in the browser \citep{prg32kit}.

\section{Why blocks, why games, why seven}

The pedagogy rests on three well-supported claims. First, children learn
powerful ideas best by building things they care about and can show to others;
this is Papert's constructionism \citep{papert1980}, and a game on a screen is
about as showable as an artefact gets. Second, block-based editors remove the
syntax errors that dominate a novice's first hours with text, without removing
the concepts; pupils who start with blocks reach comparable or better
conceptual understanding and report more confidence \citep{weintrop2017}.
Third, the habits the track practises (decomposing a problem, naming state,
testing one change at a time, reading what the machine actually did rather
than what one hoped) are the content of what \citet{wing2006} called
computational thinking, and they transfer beyond programming.

Seven is a pragmatic lower bound, not a theoretical one. The track asks a
child to read short labels, to type small numbers, to reason about ``bigger''
and ``smaller'', and to add within a few hundred. Negative numbers appear
(``change by $-3$'') and are best introduced as ``going back'' on a number
line taped to the floor. Multiplication appears once, in the score bar, and
can be treated as a recipe by younger pupils.

\begin{center}
\begin{tabularx}{\linewidth}{@{}lXX@{}}
\toprule
Age & Typical pace & Adaptations \\
\midrule
7--8 & one mission per two sessions; pairs & read aloud; pre-built starter
  projects; floor number line; teacher types titles \\
9--10 & one mission per session & pupils read the chapter; ask for a written
  prediction before each Play \\
11--13 & Missions 1--3 in two sessions & add the generated-C discussion and
  the Advanced~C editor; publish to the Store \\
14+ & one session as an on-ramp & move to Part~III (C) and compare with the
  generated source \\
\bottomrule
\end{tabularx}
\captionof{table}{Suggested pacing by age band.}
\label{tab:kitages}
\end{center}

\section{Installing the Kit for a class}

The Kit needs Python~3.11 or newer and a modern browser. The simplest
arrangement is one teacher computer running the server and every pupil device
opening its address over the classroom network; nothing is installed on pupil
devices, and tablets work because the simulator has touch controls.

\lstinputlisting[style=bash,
  caption={Running the Construction Kit with Python.}
]{scripts/construction_kit/run_python.sh}

For a computer room that is used every week, a container is easier to keep
alive. The supplied Compose file stores all projects in a \fname{data}
directory beside it, so the class's work survives restarts.

\lstinputlisting[style=bash,
  caption={Running the Construction Kit with Docker Compose.}
]{scripts/construction_kit/run_docker.sh}

\begin{center}
\begin{tabularx}{\linewidth}{@{}llX@{}}
\toprule
Variable & Default & Purpose \\
\midrule
\code{PRG32\_KIT\_DATA} & \fname{data} & where projects, sprites, and builds
  are stored \\
\code{PRG32\_KIT\_PORT} & 5090 & HTTP port \\
\code{DATABASE\_URL} & unset & use PostgreSQL instead of the default SQLite
  file \\
\code{PRG32\_STORE\_URL} & \code{http://127.0.0.1:5080} & default Cartridge
  Store \\
\code{PRG32\_BUILD\_COMMAND} & \code{python3 -m prg32 cartridge build} & how
  cartridges are compiled \\
\code{PRG32\_CART\_RAM\_KIB} & 64 & cartridge RAM profile passed to the
  builder \\
\code{SECRET\_KEY} & development value & session secret; change it for any
  shared deployment \\
\bottomrule
\end{tabularx}
\captionof{table}{Construction Kit configuration \citep{prg32kit}.}
\end{center}

\subsection*{Real cartridges are optional}

Without the \prg toolchain the Kit still does everything in
Chapters~\ref{ch:kids1}--\ref{ch:kids2} and most of Chapter~\ref{ch:kids3}:
\kbtn{Compile Cartridge} then produces a \emph{source-only} bundle containing
the generated C, the Blocks project, and a note that a toolchain is required.
To produce real \fname{.prg32} files, start the Kit from a shell in which the
\prg build command works (Chapter~\ref{ch:firstbuild}).

\lstinputlisting[style=bash,
  caption={Starting the Kit with the cartridge toolchain available.}
]{scripts/construction_kit/enable_cartridges.sh}

\begin{prgwarn}
The Kit has no pupil accounts: everyone who can reach the address sees every
project. That is convenient in a primary classroom and inappropriate on the
open Internet. Run it on the school network only. If you enable Advanced~C
with a toolchain, remember that it compiles source typed by pupils on the
teacher's machine; the Kit validates includes and entry points, but it is not
a sandbox \citep{prg32kit}. Use it with classes you trust, on a machine that
holds nothing sensitive.
\end{prgwarn}

\section{A term plan}

The three pupil chapters contain five missions. Table~\ref{tab:kitplan}
spreads them over ten 45-minute sessions, which fits a half-term club or a
unit within a computing curriculum. Each session has the same shape: a short
unplugged or whole-class start, a guided build, free modification, and a
closing circle in which two or three pupils show what they changed.

\begin{center}
\begin{tabularx}{\linewidth}{@{}clXl@{}}
\toprule
\# & Mission & Focus and outcome & Idea \\
\midrule
1 & 1 & Robot game unplugged; tour of the Kit; a square on screen & sequence \\
2 & 1 & coordinates; draw a picture from rectangles & coordinates \\
3 & 2 & buttons move the square & condition \\
4 & 2 & the fence; Step and the variable read-out & variable, debugging \\
5 & 3 & Lemon Catcher: draw, move, fall & loop, decomposition \\
6 & 3 & touching, random, score bars & collision, randomness \\
7 & 4 & beeps on events; Music Pad & event \\
8 & 4 & Sprite Lab; pixels and pictures & representation \\
9 & 5 & own game: plan on paper, then build & design \\
10 & 5 & generated C; cartridge; class showcase and publishing & abstraction \\
\bottomrule
\end{tabularx}
\captionof{table}{A ten-session plan for the Young Makers track.}
\label{tab:kitplan}
\end{center}

\subsection*{The 45-minute session}

\begin{center}
\begin{tabularx}{\linewidth}{@{}lX@{}}
\toprule
Minutes & Activity \\
\midrule
0--5 & Hook: show the finished behaviour, or play an unplugged game \\
5--20 & Guided build: the mission steps, teacher and pupils together \\
20--35 & Make it yours: the ``Try It Yourself'' changes, in pairs \\
35--40 & Look under the hood: Step, or Generate C, on one pupil's project \\
40--45 & Circle: ``what did you change, and what happened?'' Save. \\
\bottomrule
\end{tabularx}
\captionof{table}{The shape of one session.}
\end{center}

\section{Teaching moves that work}

\paragraph{Predict, run, explain.} Before a pupil presses Play after a change,
ask what will happen. A wrong prediction that is then explained teaches more
than a right result that is merely observed.

\paragraph{Pair programming.} Two pupils at one device, one ``driver'' on the
mouse and one ``navigator'' reading the mission, swapping every ten minutes.
It halves the hardware and doubles the talk.

\paragraph{One change at a time.} This is the platform's own rule for
undergraduates (Chapter~\ref{ch:platform}), and it is just as true at eight.
When a project breaks after three simultaneous changes, undo two.

\paragraph{Use Step as a slow-motion camera.} The \kbtn{Step} button advances
one frame and the grey strip shows each variable. Project one pupil's screen
and let the class call out the next value of \code{lemon\_y} before you click.

\paragraph{Let bugs be public and normal.} Keep a ``bug of the day'' on the
board, with the name of the pupil who \emph{found} it.

\begin{center}
\begin{tabularx}{\linewidth}{@{}XX@{}}
\toprule
Misconception & A response \\
\midrule
``The computer knows what I meant.'' & The Robot game; then deliberately swap
  two blocks and run it \\
``y goes up.'' & Tape a grid to the wall with (0,0) at the top left; walk it \\
``The lemon is in the basket because I can see it.'' & Change the drawn
  basket without changing the \emph{touches} numbers; nothing is caught \\
``A variable holds all its old values.'' & A physical box and one card: a new
  card replaces the old \\
``Update runs when I press a button.'' & Step with no button held; the lemon
  still falls \\
``More blocks means a better program.'' & Celebrate the shortest project that
  meets the brief \\
\bottomrule
\end{tabularx}
\captionof{table}{Common misconceptions and how to meet them.}
\end{center}

\section{From blocks to C to assembly: the vertical curriculum}

The reason to teach young children on \prg rather than on a closed blocks
environment is that nothing has to be unlearned later. The same Lemon Catcher
exists at three altitudes in this book.

\lstinputlisting[style=c,firstline=73,lastline=82,
  caption={The start of what the Kit generates for the Lemon Catcher update blocks
  (excerpt of \fname{examples/blocks\_first\_steps/lemon\_catcher.c}).}
]{examples/blocks_first_steps/lemon_catcher.c}

A pupil who has snapped ``if button LEFT pressed'' has met the conditional; the
generated \code{if (input \& PRG32\_BTN\_LEFT)} adds the fact that buttons are
bits; and Chapter~\ref{ch:asm2} shows the same test as an \code{andi} followed
by a branch. In a school that runs all three tracks, older students can be
given the younger pupils' generated C as a code-reading exercise, or asked to
rewrite a Blocks game by hand, as Appendix~\ref{app:worked} does for Lemon
Catcher, and to explain every difference.

\begin{center}
\begin{tikzpicture}[font=\sffamily\small,node distance=5mm,
  s/.style={rounded corners=3pt,draw=prgsteel,fill=prgpanel,minimum height=10mm,minimum width=40mm,align=center},
  ar/.style={-{Stealth[length=2.2mm]},draw=prgink!70,thick}]
  \node[s,fill=prgviolet!12] (b) {Blocks (Part~IV)\\\scriptsize ages 7+};
  \node[s,fill=prgamber!14,right=12mm of b] (c) {C (Part~III)\\\scriptsize ages 14+ and university};
  \node[s,fill=prggreen!12,right=12mm of c] (a) {Assembly (Part~II)\\\scriptsize university};
  \draw[ar] (b) -- node[above,font=\scriptsize]{Generate C} (c);
  \draw[ar] (c) -- node[above,font=\scriptsize]{compile} (a);
\end{tikzpicture}
\captionof{figure}{One platform, three altitudes. Each arrow is something the
tools let a learner actually look at.}
\end{center}

\section{Assessment}

Assess the ideas, not the artefact. A beautiful game copied block for block
shows less than a plain one whose author can explain it.

\begin{center}
\begin{tabularx}{\linewidth}{@{}lXXX@{}}
\toprule
Idea & Beginning & Secure & Extending \\
\midrule
Sequence & builds from a model & explains why order matters & predicts the
  effect of reordering \\
Variables & changes a given value & names and uses own variables & uses an
  expression such as \code{score * 3} \\
Conditions & uses a button block & combines conditions for a rule & expresses
  a threshold with the collision block \\
Debugging & asks for help & uses Step and the read-out & forms and tests a
  hypothesis \\
Design & modifies a game & plans a game on paper first & reviews a peer's
  game kindly and usefully \\
\bottomrule
\end{tabularx}
\captionof{table}{A simple progression rubric.}
\label{tab:kitrubric}
\end{center}

Three low-cost sources of evidence work well: the pupil's saved project
(exported as JSON from the editor); a one-minute ``explain this block to me''
conversation; and a prediction sheet on which pupils write what they expect
before pressing Play.

\section{Inclusion, safety, and data}

The simulator's game pad works with touch, mouse, and keyboard, so pupils who
cannot use one can use another. Colour is never the only carrier of meaning in
the missions, but check any pupil-designed game for colour pairs that are hard
to tell apart. Keep sessions short and screens shared.

On data: the Kit stores project titles, an optional author field, and the
projects themselves. Ask pupils to use first names or team names. Publishing
to a Cartridge Store makes a title, summary, and author string public within
that Store; publish from a teacher-held profile, to a school-run Store, after
the review described in Chapter~\ref{ch:store}. Back up the Kit by copying its
\fname{data} directory while the server is stopped.

\section{Publishing a class's games}

The Kit's \kbtn{Publish} page stores \emph{publish profiles}: a Store address
and, optionally, a bearer token. \kbtn{Compile Cartridge} on a project creates
a bundle, and \kbtn{Publish} on that bundle sends it to the Store, where it
waits as a pending submission until an editor verifies it. In a school the
natural arrangement is that the teacher holds both the Kit's publish token and
the editor role on the school Store, so that a human decision always stands
between a child's work and its publication.

\begin{prgcheck}
You are ready to teach the track when you have: run the Kit and opened it from
a second device; imported the four reference projects from
\fname{examples/blocks\_first\_steps}; played each in the simulator; used
\kbtn{Step} with a button held; clicked \kbtn{Generate C} and matched three
blocks to three lines; and decided whether your installation will build real
cartridges or source-only bundles.
\end{prgcheck}
